\begin{exercise}{34}
Let $(V,\otimes,\bbone)$ be a monoidal category. Then the functor
$\operatorname{Hom}_V(\bbone,-)\colon V\to\mathrm{Set}$ is lax monoidal.
Is it monoidal? If not: Can you find a condition on $(V,\otimes,\bbone)$
which ensures that it is?

\begin{proof}
Give $\mathrm{Set}$ its cartesian monoidal structure and put
$H(A)=\operatorname{Hom}_V(\bbone,A)$.
Write $\lambda_A\colon\bbone\otimes A\to A$ and
$\rho_A\colon A\otimes\bbone\to A$ for the inverses of
the unit isomorphisms in Definition~1.2.30.
The canonical structure maps are
\[
 \mu_{A,B}(f,g)=(f\otimes g)\delta,
 \qquad *\longmapsto\mathrm{id}_{\bbone},
 \qquad \delta=\lambda_{\bbone}^{-1}
              =\rho_{\bbone}^{-1}.
\]
Naturality follows from functoriality of $\otimes$.
Coherence for the unit isomorphisms gives
\[
 \alpha_{\bbone,\bbone,\bbone}
 (\delta\otimes\mathrm{id})\delta
 =(\mathrm{id}\otimes\delta)\delta.
\]
Postcomposing with $f\otimes(g\otimes h)$ and using naturality
of $\alpha$ proves the associativity diagram of
Definition~1.2.31. The two unit diagrams follow from naturality
of $\lambda,\rho$, which gives
\[
 \lambda_A(\mathrm{id}\otimes f)\delta=f,
 \qquad \rho_A(f\otimes\mathrm{id})\delta=f.
\]
Thus $H$ is lax monoidal.

It need not be monoidal. For
$(V,\otimes,\bbone)=(\mathrm{Ab},\otimes_{\mathbb Z},\mathbb Z)$,
identify $H(\mathbb Z)$ with $\mathbb Z$ by evaluation at $1$.
Then $\mu_{\mathbb Z,\mathbb Z}$ is multiplication
$\mathbb Z\times\mathbb Z\to\mathbb Z$, which sends both
$(2,3)$ and $(1,6)$ to $6$, so it is not an isomorphism.

A sufficient condition is that $V$ be cartesian monoidal:
$\otimes$ is the categorical product and $\bbone$ is
terminal. Then $\mu_{A,B}$ is the bijection
$(f,g)\mapsto\langle f,g\rangle$, with inverse
$k\mapsto(\operatorname{pr}_A k,\operatorname{pr}_B k)$.
Also $H(\bbone)=\{\mathrm{id}_{\bbone}\}$, so the unit
map is a bijection. Hence $H$ is monoidal in this case.
\end{proof}
\end{exercise}
